\begin{abstract}
% \mj{\\
% \href{https://people.mpi-sws.org/~dreyer/talks/talk-plmw17icfp.pdf}{How to Write Papers and Give Talks That People Can Follow}
Reinforcement learning (RL) has become a representative post-training paradigm for large language models (LLMs), enabling strong reasoning and agentic capabilities.
However, rollout generation remains a dominant \bottleneck{} because autoregressive (AR) sampling decodes responses sequentially and a small number of long-tailed generations often determine completion time.
Speculative decoding (SD) offers a natural way to address this bottleneck, as it is a well-established technique for serving fixed LLMs that reduces latency by rapidly drafting tokens and accepting them through parallel verification while preserving the target-model distribution.
However, its practical speedups do not directly carry over to RL rollouts: (i) the evolving target policy makes any fixed drafter increasingly mismatched with the policy's output distribution; and (ii) active batch sizes shrink throughout rollout decoding, shifting decoding from compute-bound to memory-bound regimes where parallel verification can exploit underutilized compute.
Therefore, accelerating RL rollouts requires both a drafter that remains effective under long, high-temperature generations from an evolving policy and system-aware use of SD that avoids compute-bound regimes.
We present \textbf{\methodtitle{}}, a system-aware self-SD framework designed to address this gap for RL rollouts.
\methodtitle{} induces a quantized drafter from the target model (\ie{}self-speculative decoding), keeping it coupled to the evolving policy without separate drafter pretraining or online adaptation.
It further coordinates a system-aware SD toggle policy with acceptance-aware draft-length adaptation, enabling speculation only in beneficial regimes while matching the drafting budget to evolving drafter quality.
\methodtitle{} reduces rollout and end-to-end latency by up to~\speeduprollout{} and~\speedupete{}, respectively, over an accelerated AR rollout baseline, while preserving final model quality.
\footnote{Code is available at \url{https://github.com/furiosa-ai/EfficientRollout}.}
\end{abstract}

% However, its practical speedups do not directly carry over to RL rollouts: (i) algorithmically, the continuously evolving target model makes static drafters stale; (ii) system-wise, active batch sizes shrink throughout rollout decoding, changing whether verification overhead can be amortized.
% As a result, existing SD approaches designed for a fixed target model do not directly translate to RL rollouts: drafting must remain effective under long, high-temperature rollout distributions and a continuously evolving policy, while speculation must avoid compute-bound regimes.